\section{Conclusion and limitations}
\label{sec:discussion}

\textbf{Conclusion.}
\looseness -1
The experiments support a simple mechanistic picture: latent
dimensions beyond the data's effective signal dimension are not
just wasted coordinates. In the RFNN they appear as a distinct noise-dimension block between the signal and sample-specific modes, and, more importantly, they lower the feature pre-activation variance so that the sample-specific modes become slower relative to the signal modes; this ratio, not the number of excess modes, is the buffer. In trainable synthetic MLPs at fixed width the same widening suppresses recall, although the theory is stated for frozen features and fixed diffusion time. In real VAE latents,
the same force competes with VAE fidelity: small bottlenecks do not
reconstruct well enough to memorize in pixel space, intermediate
dimensions both reconstruct and memorize, and very wide latents
again reduce memorization while often hurting FID. The main
takeaway is therefore predictive rather than only descriptive: once
the VAE is expressive enough to preserve image identity, the frozen
VAE latent spectrum can be turned into a mode-absorption clock that
estimates the memorization percentage over training steps for each
choice of $\dlat$.

\textbf{Practical implication.}
\looseness -1
The best latent dimension is neither the smallest code nor the
largest code. Too small loses reconstruction fidelity; too large lengthens $\taugen$ (which grows with $\dlat$ in gradient-flow time), raises the finite-$\nsamp$ score-error floor of the null block, and eventually, as $\dlat$ approaches $\nsamp$, lets memorization return. The useful regime is near the point
where VAE quality has saturated but before the diffusion model pays a
large optimization cost. The mode-absorption predictor in
Sec.~\ref{sec:spectral-predictor} makes this tradeoff concrete:
it estimates memorization risk from a frozen encoder and a 1k latent
subset.

\textbf{Implications.}
\looseness -1
This gives latent diffusion a concrete design rule. Rather than
choosing $\dlat$ from VAE quality or final FID alone, one
can estimate a memorization-risk curve before a full training sweep.
A practitioner could first use the VAE identity proxy to rule out
under-expressive latents, then use spectral pressure to rank the
remaining widths by expected memorization onset. Checkpoint audits or
early-stopping rules can then be placed near the predicted
iso-memorization times. The same calculation separates two failure
modes that look similar in final samples: low $\dlat$ fails because
the representation discards identity, while very high $\dlat$ fails
because the diffusion problem becomes wide and slow. Thus the
predictor is useful both for model selection and for deciding when a
partially trained model is likely to cross from generalization into
sample recall.

\textbf{Limitations.}
\looseness -1
(i)~The real-data diffusion experiments hold the MLP hidden width
fixed across $\dlat$ so that increasing latent dimension is not also
a capacity sweep. This isolates the latent-buffer effect, but it is
not the only scaling regime used in practice. If hidden width, depth,
or training compute are increased with $\dlat$, high-dimensional
models may recover better FID and may also memorize more because the
extra capacity can fit sample-specific directions rather than being
spent only on the harder latent score problem. (ii)~(ii) The spectral predictor is derived only for frozen features at fixed $t$: the mode clocks and the per-sample sample-bulk eigenvalue are exact or Gaussian-equivalent statements for the RFNN, the square in $P_d(s)$ is a modelling choice, and the deployed network's first layer (GELU on $[x_t;\,\text{time embedding}]$) has a nearly $\dlat$-independent pre-activation variance, so all real-data numbers use a $\tanh$ surrogate. The one-constant clock transfers between CelebA and CIFAR-10 only up to a factor $2$, and the $q$-mechanism explains roughly half of the observed real-data slowdown in log units. (iii)~The synthetic data
manifold is a linear Gaussian
mixture, so curved manifolds and class hierarchies are tested mainly
through CelebA, with CIFAR-10 included as a compact spectral
replication in the main text and a full diagnostic replication in
the appendix.
(iv)~The real-data memorization metric is pixel-space
nearest-neighbor ratio on a diverse 1k subset; other definitions may
shift the absolute percentages. (v)~The VAE and diffusion systems
interact: low-dimensional VAEs suppress measured memorization by
blurring reconstructions, while high-dimensional VAEs can degrade
FID through optimization difficulty. Separating these effects more
cleanly is the main next ablation.
